\chapter{Resultados e Discussão}
\label{cap:resultados}

%DEGAS - Mencione que a validação seguirá o formato que você mesmo previu, e referencie a Seção onde isso está // ok

Este capítulo apresenta os resultados obtidos com a implementação do protótipo, os cenários detalhados utilizados para sua homologação e uma discussão técnica sobre o papel da solução na preparação automatizada do Barema do COLCIC. A estratégia de validação adotada neste trabalho segue o planejamento previsto e os parâmetros definidos na Seção \ref{subsec:criterios_sucesso}, confrontando as saídas geradas pelo sistema com os critérios de sucesso estabelecidos para a aplicação. Todos os resultados descritos concentram-se em evidências funcionais mensuradas em ambiente de testes local.

\section{Resultados Obtidos}

%DEGAS - A velha cróitica. Cada característica do sistema merece ser explicada com certo nível de detalhe, e cada uma em um parágrafo, de preferência. //ok

O principal resultado prático deste trabalho consiste na consolidação de uma plataforma \textit{web} funcional projetada especificamente para mediar a preparação de documentos acadêmicos regulatórios. O sistema foi estruturado como um ambiente transiente de trabalho, permitindo que o discente consolide seus dados de matrícula, selecione o barema de transição aplicável, informe as horas declaradas e anexe os comprovantes correspondentes de forma assistida.

Para gerenciar as interações com os arquivos sem sobrecarregar o servidor, o sistema utiliza uma \textit{interface} baseada em componentes reativos em \textit{JavaScript} que se comunicam de maneira assíncrona com a camada de serviços do Flask\texttrademark{}. O usuário realiza a anexação de novos certificados através de campos de \textit{upload} dinâmicos; o sistema intercepta o arquivo e o envia via \textit{Fetch API} para o \textit{backend}, onde o orquestrador o aloca em uma lista de objetos temporária na memória da sessão. Para substituir ou remover documentos, o usuário conta com ícone e botões intuitivos na \textit{interface}. A substituição e a remoção de documentos ocorrem por meio da manipulação direta dos índices dessa lista volátil, garantindo respostas visuais imediatas na interface sem interrupções ou necessidade de recarregamento total da página. 

%DEGAS - De que forma o usuário pode anexar, substituir ou remover? Que módulo usa? Uma pbbreve descrição também. //ok

Outro resultado relevante foi a estruturação do mecanismo de inspeção técnica de binários operado pela biblioteca PyMuPDF\texttrademark{}. No instante em que um certificado é posicionado em uma categoria, o servidor abre o fluxo de \textit{bytes} do arquivo diretamente na memória e executa rotinas baseadas em expressões regulares para identificar padrões textuais de carga horária e datas de realização. Esse processamento imediato alimenta o motor de regras da aplicação, permitindo que o sistema atualize os contadores visuais e aplique os tetos regulamentares na \textit{interface} de maneira instantânea.

A integração entre a extração e a pré-validação de dados demonstrou utilidade prática ao atuar como um filtro de qualidade anterior à submissão oficial. Ao cruzar as informações textuais do PDF com as declarações do estudante, a plataforma sinaliza inconsistências operacionais, tais como divergências entre o nome informado e o nome identificado no certificado, declaração de carga horária superior à carga textual encontrada no documento e possível duplicidade de certificados. Dessa forma, o sistema não substitui a análise do colegiado, mas antecipa pontos de atenção que poderiam gerar retrabalho durante a conferência institucional.

%DEGAS - Mencione a utilidade da integração que tu se refere no pa´ragrafo acima //ok

\section{Validação do Protótipo}

A validação do protótipo foi conduzida por meio de testes de conformidade normativa, simulando cenários reais enfrentados pela coordenação do curso. A estrita aderência às regras do Barema é mandatória, visto que o documento final possui valor regulatório e subsidia o deferimento da colação de grau pelo Colegiado; qualquer erro de cálculo ou categorização inadequada resulta no indeferimento do processo e em prejuízos ao cronograma acadêmico do discente. Para cobrir essas variáveis, foram projetados cenários de testes funcionais, explicados a seguir.

%DEGAS - Explique (de novo, eu sei. Mas tenha pena do Leitor) porque é necessário seguir as regras do barema. Mencione que cada cenário é explicado a seguir. //ok

Nos cenários avaliados, o sistema conseguiu identificar inconsistências objetivas antes da submissão oficial. A divergência de nome foi detectada pela comparação entre o nome informado pelo discente e o nome identificado no certificado. A carga horária solicitada foi comparada com o total extraído dos documentos anexados, permitindo apontar quando o valor declarado era superior ao comprovado. A duplicidade foi verificada pela comparação do conteúdo textual dos certificados.

Também foi testado o comportamento da pré-validação quando há indisponibilidade de serviços externos de inteligência artificial. Nesse caso, o sistema retornou para a validação básica, mantendo a análise de regras objetivas sem depender exclusivamente de uma \textit{API} externa. Esse comportamento é relevante porque preserva a utilidade do sistema mesmo quando recursos adicionais de análise textual não estão disponíveis.

Além das validações individuais, foi verificada a geração do documento final. O PDF consolidado reuniu os dados preenchidos e os certificados anexados em uma estrutura única, adequada para encaminhamento à análise oficial. Esse resultado atende à necessidade de reduzir erros de montagem e de organização dos anexos.

\section{Discussão da Solução}

Os resultados indicam que a solução proposta se diferencia de ferramentas genéricas de manipulação de PDF por incorporar regras específicas do barema do COLCIC. Enquanto ferramentas externas apenas permitem juntar ou reorganizar arquivos, o sistema desenvolvido combina a montagem do documento com verificações relacionadas ao regulamento e ao conteúdo dos certificados.

A arquitetura em camadas adotada conferiu flexibilidade e isolamento de escopo ao projeto. A separação clara entre a camada de \textit{interface}, as rotas do Flask\texttrademark{}, as funções de serviço do orquestrador e os \textit{drivers} de geração de documentos permitiu que os limites normativos fossem codificados diretamente no \textit{backend} através de um dicionário de restrições explícito. Essa organização facilitou a escrita de rotinas de testes unitários automatizados, validando isoladamente os módulos como cálculo dos pisos e tetos das atividades, extração de carga horária e pré-validação.

%DEGAS O que danado é fallback? Use \textit{} para palavras estrangeiras. //ok

Apesar dos resultados promissores %promissores //ok
, a plataforma deve ser interpretada como um instrumento de apoio à decisão e conferência preliminar, não substituindo o papel avaliador da coordenação do curso. O sistema reduz os erros formais e inconsistências matemáticas evidentes antes do envio, mas a prerrogativa pedagógica de deferimento das atividades e análise de mérito dos certificados permanece estritamente institucional.

\section{Limitações Observadas}

%COMECE ELOGIANDO OS PONTOS POSITIVOS PARA DEPOIS FALAR DAS LIMITAÇÕES. E a velha crítica: explique cada uma delas em detalhe. Aqui, talvez, não precise separar em parágrafos. Mas como tu já separou a última... veja aí. //ok

Antes de delimitar as restrições do cenário atual, cabe destacar a robustez técnica demonstrada pelo protótipo no processamento estável de fluxos assíncronos e na fidelidade do documento consolidado gerado. Contudo, foram mapeadas limitações operacionais inerentes à tecnologia empregada que devem ser documentadas.

A primeira limitação reside na dependência da qualidade estrutural do arquivo PDF enviado. A extração automática de dados textuais via PyMuPDF\texttrademark{} apresenta melhor desempenho quando o documento possui uma camada de texto acessível; certificados compostos apenas por imagens ou escaneamentos de baixa resolução podem demandar o uso de Reconhecimento Óptico de Caracteres (OCR). Embora o projeto preveja suporte a esse tipo de recurso como mecanismo complementar, sua efetividade depende da qualidade do arquivo, da configuração do ambiente e da capacidade de reconhecimento do motor utilizado. Ademais, as rotinas que utilizam chamadas de inteligência artificial externa ficam condicionadas à latência de rede e à estabilidade do provedor escolhido, reforçando a importância do mecanismo de contingência estruturado no \textit{backend}.

Por fim, destaca-se que todas as homologações foram realizadas em ambiente local controlado. Embora os cenários simulados reflitam com precisão as regras do COLCIC, a ausência de testes de usabilidade com uma amostra estatisticamente heterogênea de discentes e coordenadores impede a coleta de métricas empíricas sobre a redução real de tempo administrativo e a curva de aprendizado da \textit{interface} em produção, configurando uma oportunidade para trabalhos futuros.
